\documentclass[leqno, a4paper, 12pt]{jsarticle}
\usepackage{amssymb}
\usepackage{amsthm}
\usepackage{amsmath}
\usepackage{comment}
\usepackage{url}

\usepackage{enumitem}
%\usepackage{showkeys}


%定理環境 thmはあまり使わずpropをなるべく使う。
%主張は基本的に証明の中で使い、そのため番号を付けない。
%注意環境を付けてみた。
\newtheorem{thm}{定理}
\newtheorem{prop}[thm]{命題}
\newtheorem{cor}[thm]{系}
\newtheorem{lem}[thm]{補題}
\newtheorem{df}[thm]{定義}
\newtheorem{exam}[thm]{例}
\newtheorem{fact}[thm]{事実}
\newtheorem*{claim}{主張}
\newtheorem*{cau}{注意}
\newtheorem*{rmk}{注意}
\renewcommand{\proofname}{証明}
%矢印たち。
%\toは\rightarrowと同じ。\getsは\leftarrowと同じ。同値は\iff
\newcommand{\To}{\Longrightarrow}
\newcommand{\Gets}{\LongLeftarrow}
%黒板太字
\newcommand{\nn}{\mathbb{N}}
\newcommand{\zz}{\mathbb{Z}}
\newcommand{\qq}{\mathbb{Q}}
\newcommand{\rr}{\mathbb{R}}
\newcommand{\cc}{\mathbb{C}}
%無限大
\newcommand{\mgn}{\infty}
%差集合は\setminusで統一する。等号の否定は\neq
%真部分集合は\subsetneqq　偏微分は\partial
%ディスプレイスタイル。\dfracでディスプレイ分数
\newcommand{\dpst}{\displaystyle}

\title{連結ハウスドルフ可算空間
\large{\\--- 月の裏側 ---}}
\author{ナンブ　キトラ}
\date{\today}
\begin{document}
\maketitle

月の裏側にはサンタクロースがいるらしい．
見たことはないんだが信じたことはある．

\textbf{ウリゾーンの補題の神話：}
ウリゾーンの補題というものは位相空間が正規ならば
連続関数が豊富に存在するという定理であり，さらには特徴づけにもなっているので
正規空間の性質を十分に反映している．
ウリゾーンはこの補題を連続体 ---この連続体という用語は古語であり，現代風にいうとコンパクト連結距離空間である---
に対して始めて証明しその後正規空間に対して証明した． 
それではこの補題はなんの定理のための補題であったのだろうか．補題は何かの定理のための準備であり，それはまるで太陽の光を反射して輝く月である．
太陽とは，連続体，すなわち連結コンパクト距離空間の濃度の決定であった．ウリゾーンは
退化していない(2点以上含む)連続体の濃度が
$2^{\aleph_{0}}$となることを証明するために
ウリゾーンの補題を証明したのであった．
そのような連続体を$X$とすると2点以上含むのでウリソーンの補題から
$X$上の連続関数$f\colon X\to [0, 1]$が存在して
$\{0, 1\}\subset  f(X)$となるものが存在することに注意しよう．すると$f(X)$は$[0, 1]$の連結部分集合で然るべきなので
$f(X)=[0, 1]$となり, $t\in [0, 1]$に対して
$f(p_{t})=t$となる$p_{t}$をとってくると
$X$の濃度が下から$2^{\aleph_{0}}$
で抑えられることがわかる．
上からの評価はまあ初等的であるから
$X$の濃度が決定されるというわけである．
ここでコンパクト性は$X$が正規であることを導出することのみに使われているので，上の証明においてコンパクト性を
正規性に変えても下からの濃度の評価は変わらない．
さて，正規空間の上の関数の豊富さと，それを利用しての
連続体の濃度の決定をしたウリゾーンは
この定理の「裏側」はどうなっているのか疑問に持ったのだろう．
連結コンパクトな距離空間は必ず連続体濃度をもってしまう，それでは可算な連結ハスドルフ空間は存在するのか?
という疑問にウリゾーンはたどり着いた．
ここで補足をしておくと連結可算な
\textbf{正則}空間は存在しないことに注意しよう．なぜならそのような空間は正規空間になり(一般にリンデレフかつ正則なら正規である)，
上で述べた議論から
その空間は連続体濃度以上の濃度持つからである．

\textbf{月の裏側：}
そして実際にウリゾーンはそのような連結で可算な空間を構成したのである．
ここで連結ハウスドルフ可算空間$X$はウリゾーンの補題の正反対の性質を満たすことに注意しよう．つまり連続関数
$f\colon X\to \rr$は常に定数である．このことは
$f(X)$が$\rr$の中で連結でありさらに可算集合であることからわかる．$\rr$の部分集合でそのようなものは一点集合しかない．よって連結ハウスドルフ空間は
実連続関数が全然豊富ではない空間になっている．

この文章では連結ハウスドルフで可算な空間の構成に関する論文を紹介していく．

\section{本文}

全部を読んだわけでもないし，
読んだとしても詳細に読んだわけではないが，ささやかながら
連結ハウスドルフ可算空間に関する論文を紹介していく．

 
 \subsection{概要}

連結ハウスドルフ可算空間は，ドイツ語が読めなくて，
この文章の筆者は読んでいないが
\cite{MR1512258}
によって初めて構成されたらしい．


参考文献の中にある論文の中で一番簡単な構成は
\cite{MR0060806}だと思う．
手短に連結ハウスドルフ可算空間の存在を実感したいのであれば，
これを読むと良い．
これはまた，
いわゆる1ページ論文でもある．
Bingは上半平面の可算稠密集合におおよそ
$x$軸上にあるへんな近傍系を定義して
連結ハウスドルフ可算空間を構成した．


位相空間がUrysohn空間であるとは，異なる二点に対し各々の交わらないような閉近傍が存在することを言う．
論文
\cite{MR0017527}ではウリゾーンが構成したものより条件を強めて連結なUrysohn空間でかつその上の実連続関数が常に定数となるような可算空間を構成している．

 
\subsection{整数上の位相}
次に整数の上に連結ハウスドルフ可算空間を構成した論文を紹介したい．
その前にFurstenberg
の素数の無限性の証明というものに言及せざるを得ない．
Furstenberg 
\cite{MR0068566}
は整数の上に等差数列を使った位相を導入し，
その位相空間$F$の性質を使って素数の無限性を証明した．
このFurstenbergの空間$F$は$0$次元なので全然連結ではないことに注意しよう．また$F$は距離化可能空間であることにも注意しよう．有理数の集合の位相的特徴づけを用いると$F$は$\qq$と同相であることがわかる．全然連結ではない．論文
\cite{MR3372067}ではこの空間$F$の上の具体的な距離を導入している．その距離は言うなれば``$n!$-進距離''のようなものでちょっと面白い
(多分\cite{lovas2010exotic}
はこの論文のプレプリント)．
さてそれを踏まえた上で，
Furstenbergにより等差数列を使うと$\zz$
常に変な位相が導入できることがわかったのだが，その方向性で研究をした人たちがいた．
Brown \cite{Brown}と
 Golomb,
\cite{MR0107622}による構成も鮮やかで，
整数の上に件の条件を満たす位相を構成している．
論文
\cite{MR0154249}
や
\cite{MR3905659}も参照のこと．
連立合同方程式に関する知識があればその構成はよく理解できると思う．
 Kirch
\cite{MR0239563}は彼らの構成をもう一歩推し進めて局所連結性をも充たす連結ハウスドルフ位相を整数の上に構成している．
おそらくこの整数の上の奇妙な諸位相に関する話題は
もはや一大分野を築けるくらいに膨れ上がっているのかも
知れない．


\subsection{コホモロジー論}

一般にパラコンパクトハウスドルフ空間に関しては
$H^{1}(X; \zz)\approx [X, S^{1}]$
が成り立つことが知られている．
ここで$H^{1}(X; \zz)$は$X$の一次コホモロジーで，
$[X, S^{1}]$は$X$から円周$S^{1}$
への連続写像のホモトピー同値類からなる群である．
論文
\cite{MR0224086}
では，上で述べた同型にはパラコンパクトという条件が外せないことを示している．
論文\cite{MR0224086}の本文ではないけれども，
リマークにおいて連結ハウスドルフ可算空間$Y$から輪っかを作るような構成で$Y^{\prime}$
という別の連結ハウスドルフ可算空間を構成してそれが
$H^{1}(Y^{\prime}; \zz)\not\approx [Y^{\prime}, S^{1}]$
を満たすことを述べている．
また論文\cite{MR0367974}
では上で構成された$Y^{\prime}$
の詳細な分析を行なっているらしいが，
どうやらスペイン語で書かれているらしく，
しかも
掲載
雑誌も手に入りづらい感じなので，
この文章の筆者は未確認である．

\subsection{追記}
他の論文については詳細に読んでないので省略する．
この参考文献以外にもなぜだか知らないが
連結ハウスドルフ可算空間に関する
論文はたくさんあることに留意されたい．
%READ!
%myplain is the same to amsplain style. 
\nocite{*}
\bibliographystyle{myplaindoidoi}
\bibliography{conn.bib}



\end{document}